% Options for packages loaded elsewhere
\PassOptionsToPackage{unicode}{hyperref}
\PassOptionsToPackage{hyphens}{url}
\PassOptionsToPackage{dvipsnames,svgnames,x11names}{xcolor}
%
\documentclass[
  10pt,
  a4paper,
]{article}
\usepackage{amsmath,amssymb}
\usepackage{iftex}
\ifPDFTeX
  \usepackage[T1]{fontenc}
  \usepackage[utf8]{inputenc}
  \usepackage{textcomp} % provide euro and other symbols
\else % if luatex or xetex
  \usepackage{unicode-math} % this also loads fontspec
  \defaultfontfeatures{Scale=MatchLowercase}
  \defaultfontfeatures[\rmfamily]{Ligatures=TeX,Scale=1}
\fi
\usepackage{lmodern}
\ifPDFTeX\else
  % xetex/luatex font selection
\fi
% Use upquote if available, for straight quotes in verbatim environments
\IfFileExists{upquote.sty}{\usepackage{upquote}}{}
\IfFileExists{microtype.sty}{% use microtype if available
  \usepackage[]{microtype}
  \UseMicrotypeSet[protrusion]{basicmath} % disable protrusion for tt fonts
}{}
\makeatletter
\@ifundefined{KOMAClassName}{% if non-KOMA class
  \IfFileExists{parskip.sty}{%
    \usepackage{parskip}
  }{% else
    \setlength{\parindent}{0pt}
    \setlength{\parskip}{6pt plus 2pt minus 1pt}}
}{% if KOMA class
  \KOMAoptions{parskip=half}}
\makeatother
\usepackage{xcolor}
\usepackage[margin=23mm]{geometry}
\usepackage{color}
\usepackage{fancyvrb}
\newcommand{\VerbBar}{|}
\newcommand{\VERB}{\Verb[commandchars=\\\{\}]}
\DefineVerbatimEnvironment{Highlighting}{Verbatim}{commandchars=\\\{\}}
% Add ',fontsize=\small' for more characters per line
\newenvironment{Shaded}{}{}
\newcommand{\AlertTok}[1]{\textcolor[rgb]{1.00,0.00,0.00}{\textbf{#1}}}
\newcommand{\AnnotationTok}[1]{\textcolor[rgb]{0.38,0.63,0.69}{\textbf{\textit{#1}}}}
\newcommand{\AttributeTok}[1]{\textcolor[rgb]{0.49,0.56,0.16}{#1}}
\newcommand{\BaseNTok}[1]{\textcolor[rgb]{0.25,0.63,0.44}{#1}}
\newcommand{\BuiltInTok}[1]{\textcolor[rgb]{0.00,0.50,0.00}{#1}}
\newcommand{\CharTok}[1]{\textcolor[rgb]{0.25,0.44,0.63}{#1}}
\newcommand{\CommentTok}[1]{\textcolor[rgb]{0.38,0.63,0.69}{\textit{#1}}}
\newcommand{\CommentVarTok}[1]{\textcolor[rgb]{0.38,0.63,0.69}{\textbf{\textit{#1}}}}
\newcommand{\ConstantTok}[1]{\textcolor[rgb]{0.53,0.00,0.00}{#1}}
\newcommand{\ControlFlowTok}[1]{\textcolor[rgb]{0.00,0.44,0.13}{\textbf{#1}}}
\newcommand{\DataTypeTok}[1]{\textcolor[rgb]{0.56,0.13,0.00}{#1}}
\newcommand{\DecValTok}[1]{\textcolor[rgb]{0.25,0.63,0.44}{#1}}
\newcommand{\DocumentationTok}[1]{\textcolor[rgb]{0.73,0.13,0.13}{\textit{#1}}}
\newcommand{\ErrorTok}[1]{\textcolor[rgb]{1.00,0.00,0.00}{\textbf{#1}}}
\newcommand{\ExtensionTok}[1]{#1}
\newcommand{\FloatTok}[1]{\textcolor[rgb]{0.25,0.63,0.44}{#1}}
\newcommand{\FunctionTok}[1]{\textcolor[rgb]{0.02,0.16,0.49}{#1}}
\newcommand{\ImportTok}[1]{\textcolor[rgb]{0.00,0.50,0.00}{\textbf{#1}}}
\newcommand{\InformationTok}[1]{\textcolor[rgb]{0.38,0.63,0.69}{\textbf{\textit{#1}}}}
\newcommand{\KeywordTok}[1]{\textcolor[rgb]{0.00,0.44,0.13}{\textbf{#1}}}
\newcommand{\NormalTok}[1]{#1}
\newcommand{\OperatorTok}[1]{\textcolor[rgb]{0.40,0.40,0.40}{#1}}
\newcommand{\OtherTok}[1]{\textcolor[rgb]{0.00,0.44,0.13}{#1}}
\newcommand{\PreprocessorTok}[1]{\textcolor[rgb]{0.74,0.48,0.00}{#1}}
\newcommand{\RegionMarkerTok}[1]{#1}
\newcommand{\SpecialCharTok}[1]{\textcolor[rgb]{0.25,0.44,0.63}{#1}}
\newcommand{\SpecialStringTok}[1]{\textcolor[rgb]{0.73,0.40,0.53}{#1}}
\newcommand{\StringTok}[1]{\textcolor[rgb]{0.25,0.44,0.63}{#1}}
\newcommand{\VariableTok}[1]{\textcolor[rgb]{0.10,0.09,0.49}{#1}}
\newcommand{\VerbatimStringTok}[1]{\textcolor[rgb]{0.25,0.44,0.63}{#1}}
\newcommand{\WarningTok}[1]{\textcolor[rgb]{0.38,0.63,0.69}{\textbf{\textit{#1}}}}
\setlength{\emergencystretch}{3em} % prevent overfull lines
\providecommand{\tightlist}{%
  \setlength{\itemsep}{0pt}\setlength{\parskip}{0pt}}
\setcounter{secnumdepth}{-\maxdimen} % remove section numbering
\ifLuaTeX
  \usepackage{selnolig}  % disable illegal ligatures
\fi
\usepackage{bookmark}
\IfFileExists{xurl.sty}{\usepackage{xurl}}{} % add URL line breaks if available
\urlstyle{same}
\hypersetup{
  colorlinks=true,
  linkcolor={black},
  filecolor={Maroon},
  citecolor={Blue},
  urlcolor={blue},
  pdfcreator={LaTeX via pandoc}}

\author{}
\date{}

\begin{document}

\section{DDTA BA5-T2 canonical semantic registry coverage and
governed-extension pressure test -
R1}\label{ddta-ba5-t2-canonical-semantic-registry-coverage-and-governed-extension-pressure-test---r1}

\textbf{Status:} COMPLETED / PROVISIONAL PASS WITH DOMAIN-SCOPED
REGISTRY AND GOVERNED-EXTENSION REFINEMENT

\textbf{Repository baseline tested:}
\texttt{85622414b2ff52d58f3cd11776fef4b3753afc7d}

\textbf{Scope:} execute only BA5-T2 from Work Plan R20. Do not execute
BA5 closure, BA6, AnalysisRecord/Common Finding, formal STRIDE/STRIDE-AI
schema or ThreatForge implementation design.

\subsection{1. Trial question}\label{trial-question}

Can the exact-token/no-synonym discipline established for referent names
in BA5-T1 cover the wider operative BA2 vocabulary and a realistic
extension workflow without semantic loss or authority leakage?

The falsification target is stronger than simple lookup. T2 attempts to
break the approach with:

\begin{itemize}
\tightlist
\item
  all thirteen BA2 operators;
\item
  operator-scoped role contracts;
\item
  semantic-kind keys;
\item
  controlled values;
\item
  colliding token strings across domains;
\item
  aliases;
\item
  genuinely new vocabulary;
\item
  method-local terminology; and
\item
  tool-assisted authoring.
\end{itemize}

\subsection{2. Fixed inputs}\label{fixed-inputs}

T2 treats the following as closed dependencies rather than redesign
targets:

\begin{Shaded}
\begin{Highlighting}[]
\NormalTok{BA1: BAReferent + BAProposition}

\NormalTok{BA2 proposition:}
\NormalTok{  semanticOperatorKey}
\NormalTok{  participation(roleKey, term)}
\NormalTok{  polarity}
\NormalTok{  scopedModifier(condition | temporalScope)}

\NormalTok{BA3:}
\NormalTok{  provenance / derivation / continuity}

\NormalTok{BA4:}
\NormalTok{  projection trace + shared rendering}
\NormalTok{  method{-}owned interpretation remains downstream}

\NormalTok{BA5{-}T1:}
\NormalTok{  one canonical referent name per baseline/naming scope}
\NormalTok{  no normative referent aliases}
\end{Highlighting}
\end{Shaded}

\subsection{3. Test A - thirteen operator
keys}\label{test-a---thirteen-operator-keys}

The complete current BA2 operator set is replayed under exact-token
validation:

\begin{Shaded}
\begin{Highlighting}[]
\NormalTok{transfer}
\NormalTok{produce}
\NormalTok{create}
\NormalTok{observe}
\NormalTok{transition}
\NormalTok{correlate}
\NormalTok{reference}
\NormalTok{dependOn}
\NormalTok{consumeService}
\NormalTok{realize}
\NormalTok{assignResponsibility}
\NormalTok{constrain}
\NormalTok{classify}
\end{Highlighting}
\end{Shaded}

For each operator, the authoring contract requires the registered exact
semantic key in the machine-significant binding.

Representative mutations:

\begin{Shaded}
\begin{Highlighting}[]
\NormalTok{transfer \textless{}{-} send          REJECT alias}
\NormalTok{observe  \textless{}{-} read          REJECT alias}
\NormalTok{consumeService \textless{}{-} use     REJECT alias}
\NormalTok{realize  \textless{}{-} implement     REJECT alias}
\end{Highlighting}
\end{Shaded}

This does not forbid ordinary explanatory prose. It forbids a second
normative semantic token from bypassing the registry.

\textbf{Result: PASS.} No synonym support is required.

\subsection{4. Test B - operator-scoped role
coverage}\label{test-b---operator-scoped-role-coverage}

Every closed BA2 role contract is replayed. The decisive negative
control removes the operator context.

Invalid abstraction:

\begin{Shaded}
\begin{Highlighting}[]
\NormalTok{RoleRegistry:}
\NormalTok{  actor}
\NormalTok{  source}
\NormalTok{  destination}
\NormalTok{  result}
\NormalTok{  ...}
\end{Highlighting}
\end{Shaded}

Why invalid: BA2 role validity/cardinality is operator-scoped. A
context-free row named \texttt{actor} does not tell the consumer whether
it belongs to \texttt{produce}, \texttt{create}, \texttt{observe} or
\texttt{transition}, and may hide different optionality/cardinality.

Accepted resolution:

\begin{Shaded}
\begin{Highlighting}[]
\NormalTok{OperatorRoleRegistry entry key:}
\NormalTok{  (semanticOperatorKey, roleKey)}
\end{Highlighting}
\end{Shaded}

Representative examples:

\begin{Shaded}
\begin{Highlighting}[]
\NormalTok{(transfer, source)}
\NormalTok{(transfer, destination)}
\NormalTok{(consumeService, provider)}
\NormalTok{(assignResponsibility, responsibilityKind)}
\NormalTok{(classify, semanticKind)}
\end{Highlighting}
\end{Shaded}

Mutation:

\begin{Shaded}
\begin{Highlighting}[]
\NormalTok{(transfer, src)}
\end{Highlighting}
\end{Shaded}

is rejected rather than normalized to \texttt{(transfer,\ source)}.

\textbf{Result: PASS WITH REFINEMENT.} Canonical role enforcement must
preserve operator context.

\subsection{5. Test C - cross-domain token
collision}\label{test-c---cross-domain-token-collision}

A flat global namespace is deliberately attacked:

\begin{Shaded}
\begin{Highlighting}[]
\NormalTok{project referent canonicalName = source}
\NormalTok{BA2 transfer roleKey           = source}
\end{Highlighting}
\end{Shaded}

If one global uniqueness constraint were used, this would be an
artificial collision. Semantically it is not a collision: one token
names a project referent, the other names a role in an operator
contract.

Therefore:

\begin{Shaded}
\begin{Highlighting}[]
\NormalTok{canonical uniqueness != repository{-}wide string uniqueness}
\NormalTok{canonical uniqueness = within the relevant registry domain/scope}
\end{Highlighting}
\end{Shaded}

\textbf{Result: PASS WITH REFINEMENT.} Flat global token namespace
rejected.

\subsection{6. Test D - semantic-kind
coverage}\label{test-d---semantic-kind-coverage}

T2 replays BA2 \texttt{classify}:

\begin{Shaded}
\begin{Highlighting}[]
\NormalTok{classify}
\NormalTok{  classifiedReferent {-}\textgreater{} \textless{}BAReferent\textgreater{}}
\NormalTok{  semanticKind       {-}\textgreater{} \textless{}controlled method{-}neutral key\textgreater{}}
\end{Highlighting}
\end{Shaded}

Current evidence clearly supports exact tokens such as:

\begin{Shaded}
\begin{Highlighting}[]
\NormalTok{store}
\NormalTok{contract}
\end{Highlighting}
\end{Shaded}

Mutation 1:

\begin{Shaded}
\begin{Highlighting}[]
\NormalTok{repository}
\end{Highlighting}
\end{Shaded}

If its proposed normative meaning is the same as \texttt{store}, it is
rejected as a synonym rather than admitted as a second key.

Mutation 2:

\begin{Shaded}
\begin{Highlighting}[]
\NormalTok{channel}
\end{Highlighting}
\end{Shaded}

The term is familiar and was considered during BA1 category pressure,
but T2 has no material evidence that current shared semantics cannot be
represented without a distinct \texttt{channel} kind. It is therefore
not promoted merely for convenience.

\textbf{Result: PASS.} \texttt{semanticKind} is an extensible canonical
domain; \texttt{channel} remains DEFERRED / EVIDENCE-GATED.

\subsection{7. Test E - controlled non-identifiable
values}\label{test-e---controlled-non-identifiable-values}

T2 tests whether controlled values can follow the same exact-token rule
without turning every literal into vocabulary.

Closed BA2 examples:

\begin{Shaded}
\begin{Highlighting}[]
\NormalTok{polarity:}
\NormalTok{  affirmative}
\NormalTok{  negative}

\NormalTok{scopedModifierKind:}
\NormalTok{  condition}
\NormalTok{  temporalScope}
\end{Highlighting}
\end{Shaded}

Responsibility semantics require a controlled
\texttt{responsibilityKind} domain; current evidence includes:

\begin{Shaded}
\begin{Highlighting}[]
\NormalTok{ownership}
\NormalTok{management}
\end{Highlighting}
\end{Shaded}

Mutation:

\begin{Shaded}
\begin{Highlighting}[]
\NormalTok{owner}
\end{Highlighting}
\end{Shaded}

is not silently accepted as \texttt{ownership}.

Counter-control:

\begin{Shaded}
\begin{Highlighting}[]
\NormalTok{constraintValue = 30s}
\end{Highlighting}
\end{Shaded}

A literal duration does not need a semantic registry entry merely
because it is used in a proposition. Registry admission is for reusable
controlled semantic distinctions, not every scalar/string value.

\textbf{Result: PASS WITH LOWER-BOUND REFINEMENT.} Domain-scoped
controlled values are registries; literals are not registries by
default.

\subsection{8. Test F - registry revision
mutability}\label{test-f---registry-revision-mutability}

Suppose the token \texttt{store} remains unchanged while its normative
definition is edited in place.

Then:

\begin{Shaded}
\begin{Highlighting}[]
\NormalTok{same BA baseline + token store}
\end{Highlighting}
\end{Shaded}

would not be enough to reconstruct which meaning a consumer used.

The mutation fails reproducibility.

Accepted rule:

\begin{Shaded}
\begin{Highlighting}[]
\NormalTok{registry definitions/contracts}
\NormalTok{  {-}\textgreater{} immutable revision identity}
\end{Highlighting}
\end{Shaded}

A new registry revision may add an entry while preserving old entries,
but old BA/materializations can still resolve the revision they
consumed.

The T1 referent-name case already obtains immutability from the governed
documentation baseline/naming scope.

\textbf{Result: PASS WITH REFINEMENT.} Immutable registry revision
resolution is required.

\subsection{9. Test G - genuinely new canonical
term}\label{test-g---genuinely-new-canonical-term}

T2 distinguishes five possible outcomes for an unregistered token.

\subsubsection{G1 - alias}\label{g1---alias}

\begin{Shaded}
\begin{Highlighting}[]
\NormalTok{send \textasciitilde{}= transfer}
\end{Highlighting}
\end{Shaded}

Reject new token; use \texttt{transfer}.

\subsubsection{G2 - method-specific
term}\label{g2---method-specific-term}

\begin{Shaded}
\begin{Highlighting}[]
\NormalTok{TrustBoundary}
\end{Highlighting}
\end{Shaded}

Keep it in the method/projection vocabulary. Do not promote to BA
registry.

\subsubsection{G3 - independently identifiable project
meaning}\label{g3---independently-identifiable-project-meaning}

If the supposed ``value'' needs reuse/provenance/change identity, it
becomes a \texttt{BAReferent}; it is not hidden in a controlled-value
registry.

\subsubsection{G4 - genuinely new entry in an extensible
domain}\label{g4---genuinely-new-entry-in-an-extensible-domain}

For example, a future corpus may force a new method-neutral
\texttt{semanticKind} or \texttt{responsibilityKind} value. It requires
an evidence-backed extension packet, governed acceptance and a new
immutable registry revision.

\subsubsection{G5 - new closed-contract semantic
construct}\label{g5---new-closed-contract-semantic-construct}

A genuinely new operator or role contract is not merely lexical. It
changes BA2 semantics and therefore routes to the smallest BA2 reopen.

\textbf{Result: PASS.} The extension workflow prevents BA5 from
bypassing BA1/BA2 boundaries.

\subsection{10. Test H - exact path documentation -\textgreater{} BA
-\textgreater{}
projection}\label{test-h---exact-path-documentation---ba---projection}

Representative shared binding:

\begin{Shaded}
\begin{Highlighting}[]
\NormalTok{canonical referent: cameraIngresso}
\NormalTok{operator: transfer}
\NormalTok{roles:}
\NormalTok{  source      {-}\textgreater{} cameraIngresso}
\NormalTok{  destination {-}\textgreater{} recognitionProcessor}
\NormalTok{  content     {-}\textgreater{} recognitionCapture}
\end{Highlighting}
\end{Shaded}

A shared view retains the canonical referent names and inspectable BA
keys.

A method projection may instead expose:

\begin{Shaded}
\begin{Highlighting}[]
\NormalTok{projectionOwnedKind = DataFlowEndpoint}
\NormalTok{label = cameraIngresso}
\NormalTok{trace {-}\textgreater{} BA transfer/source}
\end{Highlighting}
\end{Shaded}

or create a genuinely method-owned aggregate/interpretation with its own
label.

What it may not do is treat \texttt{sender} as an implicit BA
\texttt{source} role or make a method category a new shared
\texttt{semanticKind} without governance.

\textbf{Result: PASS.} BA4 freedom and BA5 strict canonical bindings
coexist.

\subsection{11. Test I - tool completion without semantic
authority}\label{test-i---tool-completion-without-semantic-authority}

A tool receives:

\begin{Shaded}
\begin{Highlighting}[]
\NormalTok{semanticOperatorKey = transfer}
\end{Highlighting}
\end{Shaded}

It may deterministically offer:

\begin{Shaded}
\begin{Highlighting}[]
\NormalTok{source}
\NormalTok{destination}
\NormalTok{content}
\end{Highlighting}
\end{Shaded}

from the accepted operator-role registry revision.

It may reject:

\begin{Shaded}
\begin{Highlighting}[]
\NormalTok{src}
\NormalTok{recipient}
\NormalTok{payload}
\end{Highlighting}
\end{Shaded}

unless those tokens are registered in the relevant domain (they are not
current BA2 role keys).

For an unknown \texttt{semanticKind}, it may show existing definitions
and create a candidate extension request. It cannot decide that the
unknown word is equivalent to an existing key or approve a new meaning.

\textbf{Result: PASS.} Tooling remains enforcement/assistance, not
semantic authority.

\subsection{12. Test J - authority and extension
scope}\label{test-j---authority-and-extension-scope}

T2 attacks the idea that every registry has the same lifecycle.

This collapses important distinctions:

\begin{itemize}
\tightlist
\item
  referent names are governed with the project baseline/naming scope;
\item
  BA2 operators/roles/polarity/modifier kinds belong to the methodology
  contract;
\item
  semantic-kind and other controlled-value registries may be extensible
  under their accepted method-neutral contract;
\item
  projection-local terms belong downstream.
\end{itemize}

A single mutable project registry cannot legitimately redefine
\texttt{transfer} or add a new BA2 role.

\textbf{Result: PASS WITH REFINEMENT.} Common validation discipline does
not imply common authority/lifecycle.

\subsection{13. Alias/synonym necessity
check}\label{aliassynonym-necessity-check}

T2 explicitly tries to force synonyms through:

\begin{Shaded}
\begin{Highlighting}[]
\NormalTok{send / transfer}
\NormalTok{src / source}
\NormalTok{repository / store}
\NormalTok{owner / ownership}
\NormalTok{method label / shared BA key}
\end{Highlighting}
\end{Shaded}

Every case is representable by one of:

\begin{itemize}
\tightlist
\item
  use the existing canonical key;
\item
  use free explanatory prose around a canonical binding;
\item
  keep the term method-local;
\item
  request a genuinely new evidence-backed registry entry; or
\item
  reopen the correct semantic contract if the meaning is truly new.
\end{itemize}

No case requires a normative synonym table.

\textbf{Result: NO SYNONYM/ALIAS SUPPORT FORCED.}

\subsection{14. Reopen pressure}\label{reopen-pressure}

T2 asks whether the wider registry requires semantic reopen.

\begin{Shaded}
\begin{Highlighting}[]
\NormalTok{new BAE family                         NOT FORCED}
\NormalTok{new BA2 operator                       NOT FORCED}
\NormalTok{BA1 reopen                             NOT TRIGGERED}
\NormalTok{BA2 reopen                             NOT TRIGGERED}
\NormalTok{BA3 reopen                             NOT TRIGGERED}
\NormalTok{BA4 reopen                             NOT TRIGGERED}
\NormalTok{documentation metamodel reopen         NOT TRIGGERED}
\end{Highlighting}
\end{Shaded}

The fact that a future material new operator would trigger BA2 reopen is
a routing rule, not evidence that BA2 currently fails.

\subsection{15. Trial dispositions}\label{trial-dispositions}

\begin{Shaded}
\begin{Highlighting}[]
\NormalTok{13 operator exact keys                                  PASS}
\NormalTok{operator{-}scoped role exact keys                         PASS WITH REFINEMENT}
\NormalTok{flat global registry namespace                          REJECTED}
\NormalTok{registry domain/context resolution                      REQUIRED}
\NormalTok{semanticKind controlled domain                          PASS}
\NormalTok{store / contract current canonical pressure             PASS}
\NormalTok{channel                                                  DEFERRED / EVIDENCE{-}GATED}
\NormalTok{polarity exact values                                    PASS}
\NormalTok{modifier{-}kind exact values                              PASS}
\NormalTok{responsibilityKind controlled domain                    PASS}
\NormalTok{literal value auto{-}registration                         REJECTED}
\NormalTok{immutable registry revision identity                    REQUIRED}
\NormalTok{alias candidate admission                               REJECTED}
\NormalTok{method taxonomy promotion                               REJECTED}
\NormalTok{governed extensible{-}domain extension                    PASS}
\NormalTok{closed BA2 contract mutation via BA5                     REJECTED}
\NormalTok{tool exact/context completion                           PASS}
\NormalTok{tool autonomous equivalence/approval                    REJECTED}
\NormalTok{synonym machinery                                       NOT FORCED}
\NormalTok{BA1{-}BA4 reopen                                          NOT TRIGGERED}
\NormalTok{BA5                                                      STARTED / NOT CLOSED}
\end{Highlighting}
\end{Shaded}

\subsection{16. Result}\label{result}

BA5-T2 completes with:

\begin{Shaded}
\begin{Highlighting}[]
\NormalTok{COMPLETED / PROVISIONAL PASS WITH}
\NormalTok{DOMAIN{-}SCOPED REGISTRY AND GOVERNED{-}EXTENSION REFINEMENT}
\end{Highlighting}
\end{Shaded}

The controlled-authoring hypothesis survives broader coverage. The
decisive refinements are:

\begin{enumerate}
\def\labelenumi{\arabic{enumi}.}
\tightlist
\item
  canonical registries are domain-scoped rather than one flat global
  namespace;
\item
  role keys require operator context;
\item
  registry meaning must resolve through immutable revisions;
\item
  semantic-kind/controlled-value content may be evidence-gated and
  extensible under governance;
\item
  closed BA2 operator/role semantics cannot be altered through BA5
  lexical extension; and
\item
  aliases/synonyms remain unnecessary in the normative core.
\end{enumerate}

BA5 is not closed by T2 alone.

\subsection{17. Next smallest microstep}\label{next-smallest-microstep}

Execute one integrated BA5 closure review:

\begin{quote}
\textbf{BA5-T3 - canonical registry, controlled-authoring and no-synonym
closure review}
\end{quote}

It should attempt to break the combined T1/T2 contract across referent
renames, registry revisions, human/method projections, extension
routing, current corpus examples and tool enforcement. If no material
counterexample survives, close BA5 for current thesis scope and
authorize BA6.

\end{document}
